\section{引言：为什么 Transformers 重要}

本讲的主题是 \textbf{Transformers}。这类模型建立在 self-attention 之上，已经成为现代 NLP 的默认架构，也逐渐扩展到图像、蛋白质、系统优化等领域。

\begin{importantbox}{本讲核心内容}
\begin{enumerate}
    \item Transformers 为什么在 NLP 中迅速取代 RNN
    \item self-attention 的直觉、公式与向量化实现
    \item Transformer encoder-decoder 的完整结构
    \item positional encoding、multi-head attention、masked attention 的作用
    \item Transformer 的局限，以及后续的高效变体
\end{enumerate}
\end{importantbox}

\begin{figure}[H]
\centering
\includegraphics[width=0.95\textwidth]{slides-images/slide-000.jpg}
\caption{Lecture 8 标题页：Transformers\protect\footnotemark}
\end{figure}
\footnotetext{Slides 第1页。}

\subsection{Transformers 改变了什么}

过去几年里，Transformer 不只是提升了机器翻译的质量，更改变了整个 NLP 研究范式。其核心优势不是某一个小技巧，而是把序列建模的主干从 recurrence 改成了 attention。

\begin{figure}[H]
\centering
\includegraphics[width=0.95\textwidth]{slides-images/slide-006.jpg}
\caption{原始 Transformer 论文展示的完整 encoder-decoder 架构：堆叠 self-attention、feedforward、residual、layer norm\protect\footnotemark}
\end{figure}
\footnotetext{Slides 第7页。}

\begin{figure}[H]
\centering
\includegraphics[width=0.95\textwidth]{slides-images/slide-007.jpg}
\caption{Transformer 在机器翻译上的效果：不仅 BLEU 更高，而且训练效率也更好\protect\footnotemark}
\end{figure}
\footnotetext{Slides 第8页。}

\begin{figure}[H]
\centering
\includegraphics[width=0.95\textwidth]{slides-images/slide-008.jpg}
\caption{SuperGLUE 等更难任务上的表现：Transformer 进入了更通用的 NLP 评测体系\protect\footnotemark}
\end{figure}
\footnotetext{Slides 第9页。}

\begin{figure}[H]
\centering
\includegraphics[width=0.95\textwidth]{slides-images/slide-009.jpg}
\caption{Transformer 也推动了大语言模型的兴起\protect\footnotemark}
\end{figure}
\footnotetext{Slides 第10页。}

\subsection{Transformer 不只属于 NLP}

本讲也强调了一个趋势：Transformer 不是只在文本任务里有效，而是已经扩散到更广泛的领域。幻灯片中给出了蛋白质折叠、图像分类和系统优化的例子。

\begin{figure}[H]
\centering
\includegraphics[width=0.95\textwidth]{slides-images/slide-010.jpg}
\caption{Transformer 结构开始扩展到 NLP 之外\protect\footnotemark}
\end{figure}
\footnotetext{Slides 第11页。}

\begin{figure}[H]
\centering
\includegraphics[width=0.95\textwidth]{slides-images/slide-011.jpg}
\caption{AlphaFold2 等蛋白质建模系统使用了 Transformer 思想\protect\footnotemark}
\end{figure}
\footnotetext{Slides 第12页。}

\begin{figure}[H]
\centering
\includegraphics[width=0.95\textwidth]{slides-images/slide-012.jpg}
\caption{Vision Transformer 说明 attention 也能用于图像分类\protect\footnotemark}
\end{figure}
\footnotetext{Slides 第13页。}

\begin{figure}[H]
\centering
\includegraphics[width=0.95\textwidth]{slides-images/slide-013.jpg}
\caption{Transformer 进一步被用于系统与编译器优化\protect\footnotemark}
\end{figure}
\footnotetext{Slides 第14页。}

\subsection{缩放律的启示}

幻灯片用 scaling laws 提出一个更大的问题：如果 Transformer 的性能会随着模型、数据和计算规模平滑上升，那么我们是否只需要继续扩大同样的架构？

\begin{figure}[H]
\centering
\includegraphics[width=0.95\textwidth]{slides-images/slide-014.jpg}
\caption{Transformer 与 scaling laws：性能随规模增长呈现平滑的幂律趋势\protect\footnotemark}
\end{figure}
\footnotetext{Slides 第15页。}

\begin{importantbox}{本讲的判断}
标题里问的是 ``Is Attention All We Need?''，但本讲给出的答案更接近 ``Attention 很重要，但还不够''。Transformer 的成功来自 attention + 位置编码 + 多头机制 + 残差 + layer norm + feedforward 的组合。
\end{importantbox}

\subsection{本章小结}

Transformer 之所以重要，是因为它在训练效率、长距离依赖建模和大规模扩展性上都比经典 RNN 更适合现代深度学习工作负载。接下来要理解的是：它到底是怎么工作的。

%% ========== 从 RNN 到 Attention ========== 

